% CORRECTED VERSION
% Changes:
% - Fixed algebra in Step 10 and provided a proper choice of the stepsize.
% - Added Lemma 3 linking the bounded‑domain assumption to a bound on the
%   initial optimality gap, and used it to eliminate the ad‑hoc definition of D.
% - Removed the unjustified variance term and the spurious Step 11.
% - Clarified the conditioning in Step 3 and added a short remark on the
%   role of Assumption 2 (unbiasedness) which is not needed for the main bound.
% - Updated Theorem 1 and its bound accordingly; all step annotations are
%   preserved for future reference.

\documentclass{article}
\usepackage{amsmath,amssymb,amsthm}
\usepackage{algorithm}
\usepackage{algpseudocode}
\usepackage{bm}
\usepackage{geometry}
\geometry{margin=1in}
\newtheorem{theorem}{Theorem}
\newtheorem{lemma}{Lemma}
\newtheorem{assumption}{Assumption}
\newcommand{\sign}{\operatorname{sign}}
\newcommand{\E}{\mathbb{E}}
\newcommand{\R}{\mathbb{R}}

\begin{document}

\section*{SignSGD (Sign of Stochastic Gradient Descent)}

\section*{Mathematical Formulation}
Let $f:\R^{d}\rightarrow\R$ be a (possibly non‑convex) differentiable loss.  
At iteration $k$ we have a stochastic gradient estimator $g_k\in\R^{d}$ such that  

\[
\E\bigl[g_k\mid x_k\bigr]=\nabla f(x_k),\qquad 
\E\bigl[\|g_k-\nabla f(x_k)\|_2^{2}\mid x_k\bigr]\le\sigma^{2}.
\]

The \textbf{sign update} is performed component‑wise:
\[
\boxed{ \; x_{k+1}=x_{k}-\alpha\,\sign\!\bigl(g_k\bigr)\;},
\]
where $\alpha>0$ is a (scalar) stepsize and $\sign(g_k)_i\in\{-1,+1\}$ denotes the sign of the $i$‑th coordinate of $g_k$.

When the algorithm is run in a distributed setting with $M$ workers we first collect the $M$ stochastic signs $\sign(g_k^{(m)})$, $m=1,\dots,M$, and replace $\sign(g_k)$ by the coordinate‑wise \emph{median} of these $M$ signs.  The convergence analysis below only requires the following probabilistic guarantee on each coordinate (the “median‑based variance bound’’).

\section*{Pseudocode}
\begin{algorithm}[H]
\caption{SignSGD with optional majority‑vote (median) aggregation}
\begin{algorithmic}[1]
\Require Initial point $x_0\in\R^{d}$, stepsize $\alpha>0$, number of iterations $T$,
        number of workers $M$ (set $M=1$ for the single‑machine case)
\For{$k=0$ \textbf{to} $T-1$}
    \State \textbf{Parallel on each worker $m$}:
    \State \hspace{1em} Sample a mini‑batch $\mathcal B_k^{(m)}$ and compute stochastic gradient $g_k^{(m)}$.
    \State \hspace{1em} Compute coordinate‑wise sign $s_k^{(m)}:=\sign\!\bigl(g_k^{(m)}\bigr)$.
    \State \textbf{Aggregation:}
    \State \hspace{1em} For each coordinate $i$, set
    \[
      \tilde s_{k,i}= \operatorname{median}\bigl\{s_{k,i}^{(1)},\dots,s_{k,i}^{(M)}\bigr\}.
    \]
    \State Update the model:
    \[
      x_{k+1}=x_{k}-\alpha\,\tilde s_{k}.
    \]
\EndFor
\State \textbf{Return} $x_T$ (or the average $\frac1T\sum_{k=0}^{T-1}x_k$).
\end{algorithmic}
\end{algorithm}

\section*{Dry Run Example}
Consider the univariate quadratic $f(w)=(w-3)^{2}$, whose gradient is $\nabla f(w)=2(w-3)$.  
Take $w_{0}=0$, stepsize $\alpha=0.1$, and a deterministic “stochastic’’ gradient $g_{k}= \nabla f(w_{k})$ (so the sign is exact).

\begin{tabular}{c|c|c|c}
Iteration $k$ & $w_{k}$ & $g_{k}=2(w_{k}-3)$ & Update $w_{k+1}=w_{k}-0.1\sign(g_{k})$ \\ \hline
0 & $0$ & $-6$ & $0-0.1(-1)=0.1$ \\
1 & $0.1$ & $-5.8$ & $0.1-0.1(-1)=0.2$ \\
2 & $0.2$ & $-5.6$ & $0.2-0.1(-1)=0.3$ \\
\end{tabular}

The objective values are $f(0)=9$, $f(0.1)=8.41$, $f(0.2)=7.84$, $f(0.3)=7.29$, showing monotone decrease.

\newpage

\section*{Assumptions}
\begin{assumption}[Smoothness]\label{ass:smooth}
$f$ is $L$‑smooth: for all $x,y\in\R^{d}$,
\[
f(y)\le f(x)+\langle\nabla f(x),y-x\rangle+\frac{L}{2}\|y-x\|_{2}^{2}.
\]
\end{assumption}

\begin{assumption}[Unbiased stochastic gradient]\label{ass:unbiased}
For every $k$,
\[
\E\!\left[g_k\mid x_k\right]=\nabla f(x_k),\qquad
\E\!\left[\|g_k-\nabla f(x_k)\|_{2}^{2}\mid x_k\right]\le\sigma^{2}.
\]
\end{assumption}
% NOTE: This assumption is not needed for the main bound; it is kept for completeness.

\begin{assumption}[Coordinate‑wise sign correctness (median bound)]\label{ass:sign}
There exists a constant $\rho\in(0,\tfrac12]$ such that for each coordinate $i\in\{1,\dots,d\}$,
\[
\Pr\!\bigl(\sign(g_{k,i})=\sign(\nabla f(x_k)_i)\mid x_k\bigr)\;\ge\;\tfrac12+\rho .
\tag{1}
\]
If a median of $M$ independent workers is used, the same bound holds with $\rho$ replaced by
\[
\rho_M=1-2\exp\!\bigl(-2M\rho^{2}\bigr),
\]
by Hoeffding’s inequality.
\end{assumption}

\begin{assumption}[Bounded domain]\label{ass:bounded}
All iterates remain in a set $\mathcal X$ of diameter $D$, i.e.
\[
\|x_k-x^{*}\|_{2}\le D\qquad\text{for some (unknown) minimizer }x^{*}.
\]
\end{assumption}

\section*{Main Convergence Theorem}
\begin{theorem}[Convergence of SignSGD]\label{thm:main}
Assume \ref{ass:smooth}–\ref{ass:bounded}. Let the constant stepsize be
\[
\alpha \;=\; \sqrt{\frac{2\,(f(x_0)-f^{*})}{L\,d\,T}}\;,
\qquad\text{where }f^{*}:=\inf_{x}f(x).
\]
Then the iterates produced by SignSGD satisfy
\[
\frac{1}{T}\sum_{k=0}^{T-1}\E\!\bigl[\|\nabla f(x_k)\|_{2}\bigr]
\;\le\;
\frac{\sqrt{2\,L\,(f(x_0)-f^{*})}}{\rho\,\sqrt{T}} .
\tag{2}
\]
In particular, for deterministic gradients ($\sigma^{2}=0$) we obtain the
rate $O\!\bigl(1/\sqrt{T}\bigr)$.
\end{theorem}

\section*{Theoretical Properties}
\begin{itemize}
\item \textbf{Stability.}  From the descent inequality (Step 4) we see that
$f(x_{k+1})\le f(x_k)$ whenever $\alpha\le\frac{2\rho}{L d}$.
\item \textbf{Hyper‑parameter sensitivity.}  The bound (2) depends on $\alpha$
only through the classic bias–variance trade‑off; the chosen $\alpha$
balances the two terms and yields the optimal $O(1/\sqrt{T})$ rate.
\item \textbf{Comparison to SGD.}  For smooth non‑convex problems SGD attains
$O(1/\sqrt{T})$ in terms of $\E\|\nabla f\|_{2}^{2}$.  SignSGD replaces the
Euclidean inner product by an $\ell_{1}$‑type inner product, incurring a
$\sqrt{d}$ factor in the intermediate bound, which disappears after the
stepsize optimisation.
\end{itemize}

\section*{Preliminaries (Definitions and Lemmas)}
\begin{lemma}[Smoothness inequality]\label{lem:smooth}
If $f$ is $L$‑smooth then for any $x$ and any vector $u$,
\[
f(x-u)\le f(x)-\langle\nabla f(x),u\rangle+\frac{L}{2}\|u\|_{2}^{2}.
\]
\end{lemma}
\begin{proof}
Apply Definition \ref{ass:smooth} with $y=x-u$ and rearrange.
\end{proof}

\begin{lemma}[Sign‑correctness expectation]\label{lem:signexp}
Under Assumption \ref{ass:sign},
\[
\E\!\bigl[\nabla f(x)_i\,\sign(g_i)\mid x\bigr]\;\ge\;2\rho\,\bigl|\nabla f(x)_i\bigr|
\quad\forall i.
\tag{3}
\]
\end{lemma}
\begin{proof}
Let $a_i:=\nabla f(x)_i$ and $s_i:=\sign(g_i)$.  By (1),
\[
\Pr(s_i=\sign(a_i))\ge\tfrac12+\rho,\qquad
\Pr(s_i=-\sign(a_i))\le\tfrac12-\rho .
\]
Hence
\[
\E[a_i s_i]
=|a_i|\bigl[\Pr(s_i=\sign(a_i))-\Pr(s_i=-\sign(a_i))\bigr]
\ge|a_i|(2\rho)=2\rho|a_i|.
\]
\end{proof}

\begin{lemma}[Bound on the initial optimality gap from the bounded domain]\label{lem:gap}
If $\|x_0-x^{*}\|_{2}\le D$ and $f$ is $L$‑smooth, then
\[
f(x_0)-f^{*}\;\le\;\frac{L}{2}\,D^{2}.
\tag{4}
\]
\end{lemma}
\begin{proof}
Apply Lemma \ref{lem:smooth} with $x=x^{*}$, $u=x_0-x^{*}$ and use $f(x^{*})=f^{*}$.
\end{proof}

\section*{Main Proof (Step‑by‑Step Derivation)}
We prove Theorem \ref{thm:main}.  Throughout let $x^{*}$ be any point attaining $f^{*}=\inf_{x}f(x)$.

\bigskip
\noindent\textbf{[Step 1]} Apply Lemma \ref{lem:smooth} with $u=\alpha\sign(g_k)$:
\[
f(x_{k+1})\le f(x_k)-\alpha\langle\nabla f(x_k),\sign(g_k)\rangle
               +\frac{L\alpha^{2}}{2}\|\sign(g_k)\|_{2}^{2}.
\tag{5}
\]

\noindent\textbf{[Step 2]} Since each component of $\sign(g_k)$ lies in $\{-1,+1\}$,
\[
\|\sign(g_k)\|_{2}^{2}=d.
\tag{6}
\]

\noindent\textbf{[Step 3]} Take conditional expectation w.r.t.\ the stochastic gradient $g_k$ given $x_k$.
Using Lemma \ref{lem:signexp} (the expectation is taken over $g_k$) and linearity,
\[
\E\!\bigl[\langle\nabla f(x_k),\sign(g_k)\rangle\mid x_k\bigr]
   =\sum_{i=1}^{d}\E\!\bigl[\nabla f_i(x_k)\sign(g_{k,i})\mid x_k\bigr]
   \ge 2\rho\sum_{i=1}^{d}|\nabla f_i(x_k)|
   =2\rho\|\nabla f(x_k)\|_{1}.
\tag{7}
\]

\noindent\textbf{[Step 4]} Insert (6)–(7) into (5) and then take the full expectation:
\[
\E\bigl[f(x_{k+1})\bigr]
\le \E\bigl[f(x_k)\bigr]
   -2\alpha\rho\,\E\!\bigl[\|\nabla f(x_k)\|_{1}\bigr]
   +\frac{L\alpha^{2}d}{2}.
\tag{8}
\]

\noindent\textbf{[Step 5]} Rearrange (8) to isolate the gradient term:
\[
2\alpha\rho\,\E\!\bigl[\|\nabla f(x_k)\|_{1}\bigr]
\le \E\bigl[f(x_k)-f(x_{k+1})\bigr]
   +\frac{L\alpha^{2}d}{2}.
\tag{9}
\]

\noindent\textbf{[Step 6]} Sum (9) over $k=0,\dots,T-1$ and telescope the function values:
\[
2\alpha\rho\sum_{k=0}^{T-1}\E\!\bigl[\|\nabla f(x_k)\|_{1}\bigr]
\le f(x_0)-\E\bigl[f(x_T)\bigr]
   +\frac{L\alpha^{2}dT}{2}
\le f(x_0)-f^{*}+\frac{L\alpha^{2}dT}{2}.
\tag{10}
\]

\noindent\textbf{[Step 7]} Divide (10) by $2\alpha\rho T$:
\[
\frac{1}{T}\sum_{k=0}^{T-1}\E\!\bigl[\|\nabla f(x_k)\|_{1}\bigr]
\le \frac{f(x_0)-f^{*}}{2\alpha\rho T}
   +\frac{L\alpha d}{4\rho } .
\tag{11}
\]

\noindent\textbf{[Step 8]} Relate the $\ell_{1}$‑norm to the $\ell_{2}$‑norm:
\[
\|\nabla f(x_k)\|_{2}\le\|\nabla f(x_k)\|_{1}
\le\sqrt{d}\,\|\nabla f(x_k)\|_{2}.
\]
Thus
\[
\frac{1}{T}\sum_{k=0}^{T-1}\E\!\bigl[\|\nabla f(x_k)\|_{2}\bigr]
\le \frac{1}{\sqrt{d}}\,
      \frac{1}{T}\sum_{k=0}^{T-1}\E\!\bigl[\|\nabla f(x_k)\|_{1}\bigr].
\tag{12}
\]

\noindent\textbf{[Step 9]} Combine (11) and (12):
\[
\frac{1}{T}\sum_{k=0}^{T-1}\E\!\bigl[\|\nabla f(x_k)\|_{2}\bigr]
\le
\frac{f(x_0)-f^{*}}{2\alpha\rho\sqrt{d}\,T}
   +\frac{L\alpha\sqrt{d}}{4\rho } .
\tag{13}
\]

\noindent\textbf{[Step 10]} \textbf{(Corrected)} Choose $\alpha$ to balance the two
terms on the right‑hand side of (13).  Setting the derivative of the RHS with
respect to $\alpha$ to zero yields
\[
\alpha^{*}
\;=\;
\sqrt{\frac{2\,(f(x_0)-f^{*})}{L\,d\,T}} .
\tag{14}
\]
Substituting (14) into (13) gives
\[
\frac{1}{T}\sum_{k=0}^{T-1}\E\!\bigl[\|\nabla f(x_k)\|_{2}\bigr]
\;\le\;
\frac{\sqrt{2\,L\,(f(x_0)-f^{*})}}{\rho\,\sqrt{T}} .
\tag{15}
\]
This is exactly the bound stated in \eqref{thm:main}.

\noindent\textbf{[Step 11]} \textbf{(Removed)} The previous version introduced an
additional variance term $\frac{\alpha\sigma^{2}}{2\rho D}$ without any
derivation.  Since the sign‑correctness assumption \ref{ass:sign} already
captures the effect of stochastic noise on the direction, no extra term
appears in the final bound; consequently the proof for the noisy case is
identical to the deterministic one.

\hfill $\square$

\section*{Rate Derivation (With Constants)}
Collecting the constants from the steps above, the final bound reads
\[
\boxed{
\frac{1}{T}\sum_{k=0}^{T-1}\E\!\bigl[\|\nabla f(x_k)\|_{2}\bigr]
\;\le\;
\frac{\sqrt{2\,L\,(f(x_0)-f^{*})}}{\rho\,\sqrt{T}} .
}
\]
Using Lemma~\ref{lem:gap} and the bounded‑domain assumption,
$f(x_0)-f^{*}\le \tfrac{L}{2}D^{2}$, we may also write
\[
\frac{1}{T}\sum_{k=0}^{T-1}\E\!\bigl[\|\nabla f(x_k)\|_{2}\bigr]
\;\le\;
\frac{L D}{\rho\,\sqrt{T}} .
\]

\section*{Special Cases}
\begin{itemize}
\item \textbf{Deterministic gradients ($\sigma=0$).}  The bound above
reduces to the pure $O(1/\sqrt{T})$ rate (15).
\item \textbf{Large number of workers $M$.}  Replacing $\rho$ by
$\rho_M=1-2e^{-2M\rho^{2}}$ (Hoeffding) improves the constant, effectively
accelerating convergence.
\item \textbf{Convex $f$.}  When $f$ is convex, the same proof yields a bound
on the suboptimality $f(\bar x_T)-f^{*}$ of order $O(1/\sqrt{T})$.
\end{itemize}

\end{document}

% ===== CORRECTION SUMMARY =====
% 1. [Step 10] Issue: algebraic error and undefined D. Fixed by explicitly solving the
%    bias–variance balance, yielding the stepsize (14) and the correct bound (15).
% 2. [Step 11] Issue: unjustified variance term. Removed the spurious step and
%    clarified that the sign‑correctness assumption already accounts for stochastic
%    noise, so no extra $\sigma^{2}$ term appears.
% 3. Missing usage of Assumption 4: added Lemma 3 (bound (4)) linking the diameter
%    $D$ to the initial optimality gap and used it to express the final bound
%    in terms of $D$ if desired.
% 4. Step 3 clarification: explicitly stated that the expectation is taken with
%    respect to $g_k$ conditioned on $x_k$.
% 5. Overall: ensured every assumption is either used or explicitly noted as
%    not required for the main result, and all step annotations are retained.